\subsection{The full primitive logarithm and its exact exponential threshold}

\begin{theorem}[DP41: optimal multivariate inflation for the primitive class]\label{DP41}
Let $p$ be odd, $h\ge2$, $d=h-1$, $q=p^h$, and $\mathcal O=W(\F_q)$. For the actual Honda line and its Cartier dual use the Teichm\"uller coordinates and invariant primitive frame of DP38. Put
\[
 s=\frac1{p-1},\qquad \rho=\frac1{q-1},\qquad
 r_0=\rho,\qquad r_j=s-p^j\rho\quad(1\le j\le d).
\]
Let $b(S,\mathbf Y)$ be the invariant rational logarithmic splitting and
$\beta(x,\mathbf y)=b(E_0(x),E_H(\mathbf y))$ its expression in additive coordinates. For $\lambda\in\mathcal O^\times$ and every rational $a\ge0$,
\[
 w_{a\mathbf r}\bigl(\exp(\lambda\beta)\bigr)=
 \begin{cases}-\infty,&a<p,\\0,&a\ge p.\end{cases}
\]
The same formula holds for $\exp(\lambda b)$ in group coordinates. Multiplication by any nonzero integral $U(E_0,E_H)$ preserves this exact boundedness threshold. Above the threshold the product is integral after contraction, with finite coefficient infimum; that infimum is zero if $U(0)$ is a unit. Thus even a nonunit integral prefactor cannot improve the threshold. Here $w_{\mathbf v}(\sum a_\alpha Z^\alpha)=\inf_\alpha(\vp(a_\alpha)+\mathbf v\cdot\alpha)$.

The entire primitive logarithm, with Zink's negative $\operatorname{hex}$ convention, is
\begin{equation}\label{eq:fullprimitive}
\begin{split}
 b(S,\mathbf Y)
 ={}&-\sum_{j=1}^d\sum_{l,n\ge0}
       \frac{S^{p^{n+lh+j}}Y_j^{p^n}}{p^{n+l+1}}\\
 &+\left(\sum_{u\ge0}\frac{S^{q^u}}{p^u}\right)
      \sum_{j=1}^d\sum_{l\ge0}
       \frac{Y_j^{p^{lh+h-j}}}{p^{dl+h-j}}.
\end{split}
\end{equation}
The logarithms in these exact coordinates are $L_0=\sum S^{q^u}/p^u$ and the $L_H$ of DP32. The threshold concerns the full multivariate series; it does not by itself determine a restriction or a selected harmonic value.
\end{theorem}

\begin{proof}
\emph{Divided ghosts and logarithms.}
Use DP38's display matrices
\[
 Ae_0=e_1,\quad Ae_j=pe_{j+1}\ (j\ge1),\qquad
 A'e_0^*=pe_1^*,\quad A'e_j^*=e_{j+1}^*\ (j\ge1),
\]
with cyclic indices. Thus $\langle Au,A'v\rangle=p\langle u,v\rangle$,
$A^h=p^d$, and $(A')^h=p$. Write $u_n(z)=w_n(z)/p^n$ for divided Witt ghosts. The rational resolution operator is $A\operatorname{shift}-1$, with inverse
$-\sum_{r\ge0}(A\operatorname{shift})^r$ on finite-support nilpotent ghosts. The additive logarithm of a lift is
\begin{equation}\label{eq:ghostlog}
 \operatorname{pr}_T\sum_{n\ge0}A^nu_n(z).
\end{equation}
On the relation module this telescopes to $-\operatorname{pr}_Tu_0(Q)=0$, and on the section with only a zeroth ghost in $T$ it is the identity.

All these operations are justified on the integral construction before inverting $p$. On a truncation $N^D=0$ with $D\ge2$, the formula
\[
 u_n(z)=\sum_{r=0}^{\min(n,C)}
              z_{n-r}^{p^r}/p^r,\qquad C=\floor{\log_p(D-1)},
\]
uniformly bounds the denominators of divided ghosts of integral Witt vectors. Since the matrices contract after each complete cycle, the sums and telescoping identities converge on the completed integral resolution. The particular Teichm\"uller sections and rational zeroth-ghost sections below have finite ghost support on each nilpotent truncation; passage to formal series is coefficientwise. This proves both completion and rationalization steps. The source context is Zink (147)--(156), pp.~204--207 \cite{Zink}.

For $\sigma(S)=[S]e_0$, the ghosts are $S^{p^n}e_0/p^n$, and cycles returning to $e_0$ in \eqref{eq:ghostlog} give exactly $L_0(S)$. For $\tau(\mathbf Y)=\sum_{j=1}^d[Y_j]e_j^*$, put
$v_n=\sum_jY_j^{p^n}e_j^*/p^n$ and
$Z=\sum_{n\ge0}(A')^nv_n$. It satisfies
$Z=\sum_jY_je_j^*+(A'/p)Z(\mathbf Y^p)$.
Write $C(\mathbf Y)$ for its $e_0^*$ component and $L_j$ for the others. Then
\[
 C=p^{-1}L_d(\mathbf Y^p),\qquad
 L_1=Y_1+p^{-1}L_d(\mathbf Y^{p^2}),\qquad
 L_j=Y_j+p^{-1}L_{j-1}(\mathbf Y^p)\ (j\ge2).
\]
These are exactly DP32's logarithm equations. The path from $e_j^*$ to $e_0^*$ has length $h-j$, and each further cycle contributes one numerator factor $p$, so
\[
 C(\mathbf Y)=\sum_{j=1}^d\sum_{l\ge0}
                \frac{Y_j^{p^{lh+h-j}}}{p^{dl+h-j}}.
\]

\emph{Comparison with an additive bilinear section.}
Let $x=L_0(S)$ and $\mathbf y=L_H(\mathbf Y)$. Choose additive rational lifts $g_0(x)$, $g'_0(\mathbf y)$ with only zeroth ghosts $xe_0$, $\sum y_je_j^*$. The scalar-one section in that additive frame is bilinear. For the comparison defects put
\[
 k=(A\operatorname{shift}-1)^{-1}(\sigma-g_0),\qquad
 l=(A'\operatorname{shift}-1)^{-1}(\tau-g'_0).
\]
Their zeroth $T$ and $T'$ projections vanish by \eqref{eq:ghostlog}, so they belong to the rationalized normal submodules. Explicitly,
\[
 k_n=-\sum_{r\ge0}A^ru_{n+r}+\mathbf1_{n=0}xe_0,\qquad
 l_n=-\sum_{r\ge0}(A')^rv_{n+r}
                 +\mathbf1_{n=0}\sum_jy_je_j^*.
\]
Zink's Proposition~106, p.~226, and (218), p.~239, express the Teichm\"uller frame as the additive frame multiplied by
\[
 \operatorname{hex}\langle V^{-1}k,\tau\rangle\,
 \operatorname{hex}\langle g_0,l\rangle^{-1}.
\]
The bilinear splitting scalar in that frame is the inverse product. Zink's (196), p.~232, and (210)--(211), p.~236, have the convention
\[
 \log\operatorname{hex}(z)=-\sum_{n\ge0}w_n(z)/p^n.
\]
For two divided-ghost sequences the pairing therefore has logarithm
$-\sum_np^n\langle u_n,v_n\rangle$. Substitution gives
\[
 b=-\sum_{n,r\ge0}p^n\langle A^{r+1}u_{n+r+1},v_n\rangle
          +L_0(S)C(\mathbf Y).
\]
A path of length $lh+j$ from $e_0$ ends at $e_j$ with numerator $p^{dl+j-1}$. Combining both ghost denominators and the factor $p^n$ yields $p^{-n-l-1}$, proving \eqref{eq:fullprimitive}. Every fixed ordinary degree has finitely many summands. Some terms may cancel, but no disjoint-support assumption was made. Every monomial is invariant under
$(\zeta S,\zeta^{-p}Y_1,\ldots,\zeta^{-p^d}Y_d)$; both axes vanish. The additive lifts are equivariant, so the uniqueness argument of DP37--DP38 for this fixed frame identifies this with its invariant primitive splitting. No bilinear correction has been inserted.

\emph{Exact initial form after contraction.}
Fix rational $a\ge p-1$ and choose powers $\alpha_i$ of a common ramified uniformizer with $\vp(\alpha_i)=ar_i$. Set
$b_a=b(\alpha_0S,\alpha_1Y_1,\ldots,\alpha_dY_d)$ and $t=as-1\ge0$.
All monomials of unscaled weight $s$ have the same scalar after division by $p$; call it $\mu$, with $\vp(\mu)=t$. We claim
\begin{equation}\label{eq:primitiveinitial}
 \vp(b_a)=t,\qquad
 \overline{\mu^{-1}b_a}
       =-\sum_{j=1}^dS^{p^j}Y_j+SY_d^p=:Q.
\end{equation}
For a term in the first family of \eqref{eq:fullprimitive}, its weight is
$p^n(s+p^j\rho(q^l-1))$. Its scaled valuation minus $t$ is
\[
 as(p^n-1)-n+a p^{n+j}\rho(q^l-1)-l.
\]
The first part is at least $p^n-1-n$, strictly positive for $n\ge1$ since $p$ is odd. The second is at least $((p-1)p^{n+j}-1)l$, strictly positive for $l\ge1$. Only $n=l=0$ survives.

For the product family put $k=h-j$. Its coefficient valuation is $-u-dl-k$ and weight
$W=\rho q^u+q^l(sp^k-q\rho)$. The baseline $u=l=0,k=1$ has weight $s$, and
\[
 W-s=\rho(q^u-1)+(q^l-1)(sp^k-q\rho)+s(p^k-p).
\]
After multiplication by $p-1$, the three terms are at least
$(p-1)u$, $(q-p)l$, and $p^k-p$, respectively. Here
$sp^k-q\rho\ge sp-q\rho=s-\rho>0$.
Consequently the scaled valuation minus $t$ is at least
\[
 (p-2)u+(q-p-d)l+(p^k-p-k+1).
\]
Every term is nonnegative and equality requires $u=l=0,k=1$.
The surviving monomials in \eqref{eq:primitiveinitial} are distinct, and
$\partial Q/\partial Y_1=-S^p\ne0$, including when $d=1$.

\emph{Exponential obstruction, including arbitrary nonzero bounded multipliers.}
We prove the needed algebraic lemma in full. Let $R'$ be a complete discrete valuation ring over $\Z_p$, with uniformizer $\pi$, and let $z$ have zero constant term and bounded coefficients. Suppose its coefficient infimum is $t\in[0,s)$ and its initial form has a nonzero formal partial derivative. Put
$B=R'[[Z_1,\ldots,Z_m]][1/\pi]$.
Its coefficient infimum is multiplicative: remove each attained minimum power of $\pi$ and reduce the product in the domain $k'[[\mathbf Z]]$. It extends to a valuation on $\operatorname{Frac}(B)$, with residue field $\operatorname{Frac}(k'[[\mathbf Z]])$.

The series $\exp(pz)$ is integral and
\[
 \vp(\exp(pz)-1)=1+t
\]
with the initial form of $pz$. Indeed every exponential order $n\ge2$ has valuation at least $n(1+t)-\vp(n!)>1+t$, since
$\vp(n!)\le(n-1)s$ and $1+t-s>0$.
Suppose a nonzero $f\in B$ made $fE$ bounded, where $E=\exp z$.
Then $E=(fE)/f$ belongs to $\operatorname{Frac}(B)$. The identity
$E^p=\exp(pz)$ gives $\vp(E)=0$ and $\bar E^p=1$, hence $\bar E=1$.
Put $\delta=\vp(E-1)>0$. If $\delta\ge s$, every term in
$(1+(E-1))^p-1$ has valuation at least $1+s>1+t$, a contradiction.
If $\delta<s$, its unique least term is $(E-1)^p$. Its initial form is a scalar multiple of a $p$th power in the residue fraction field. Formal partial derivatives extend to that field by the quotient rule and kill such powers, whereas the initial form of $pz$ has a nonzero derivative. This is again a contradiction. Thus every nonzero bounded $f$ gives unbounded coefficients for $f\exp z$.

Apply the lemma to $z=\lambda b_a$. For $p-1\le a<p$, \eqref{eq:primitiveinitial} has $0\le t<s$, so every nonzero bounded multiplier leaves coefficient infimum $-\infty$. For $a\ge p$, $t\ge s$ and
\[
 \vp((\lambda b_a)^n/n!)\ge ns-\vp(n!)
                    =s_p(n)/(p-1)\ge0.
\]
Its exponential is integral with constant one. For smaller $a$, monotonicity of weighted infima supplies the unbounded conclusion.

\emph{Coordinate changes and prefactors.}
For $a\ge1$, both contracted logarithms and their inverses are integral with identity differential. For the line this follows from
$a(q^u-1)/(q-1)-u\ge0$; for the dual it follows from DP32's path-weight increments, each at least its number of arrows. Integral substitution with an integral inverse preserves the extended coefficient infimum: the two substitution inequalities prove equality in the bounded case, and an allegedly bounded image of an unbounded series would become bounded on applying the inverse.
This transfers the result to $\beta$ for $a\ge1$; monotonicity handles $a<1$.
For a nonzero integral $U$, its expression under these contracted inverses is bounded and nonzero for $a\ge1$. The lemma handles $p-1\le a<p$, monotonicity all smaller $a$, and integrality handles $a\ge p$. If $U(0)$ is a unit, its contraction is an integral unit and the finite infimum is zero. This proves every assertion.
\end{proof}

DP41 proves that DP36's inflation factor $p$ is exact in this family of weights. The arbitrary-multiplier conclusion uses the nonzero-derivative initial form and the identity $E^p=\exp(pz)$; it does not generalize to every series with unbounded denominators, as DP35's example $f=p-z$ shows. A path can still annihilate a prefactor or cancel an initial form. The next theorem treats the primitive mixed paths explicitly.

